\chapter{Why the Text Is Written This Way}
\label{ch:method_consequence}

\epigraph{The recursive structure, explicit assumption rebuilding, and avoidance
of magical terminology jumps in this monograph are not aesthetic choices. They
are direct consequences of the epistemological commitments developed in Part~I.}{}

\section{Writing Style as Mechanistic Consequence}

A text that argues for provenance preservation and reconstruction transparency
while itself operating through aggressive compression and assumed priors would be
internally inconsistent. The monograph's form is therefore a mechanistic
consequence of its content, not a stylistic preference.

The specific consequences are:

\textbf{Progressive definition accumulation.} Terms are not assumed stable across
disciplinary boundaries. Each term is introduced when it is first needed and
derived from previously established concepts where possible. The result is a text
that restates earlier definitions at each new level of abstraction---which reads
as repetition under a compression-first reading but functions as triangulation
under a reconstruction-first reading.

\textbf{Explicit cross-referencing.} When a later argument depends on an earlier
definition, the dependency is marked. This is the textual implementation of the
provenance requirement: the dependency graph of the argument is made visible
rather than hidden.

\textbf{Appendices as provenance maps.} The appendices function not as supplementary
material but as the boot sequences of the formal frameworks used throughout.
A reader who wants to understand why a particular claim was made can follow the
reference chain back through the appendices to the primitive definitions that
ground it.

\textbf{Constraint boundary sections.} Each major chapter ends with an explicit
audit of what the chapter establishes, what it does not, what remains speculative,
and what would falsify its claims. These sections implement the monograph's
argument against opaque abstraction at the level of its own structure.

\section{Against Disciplinary Compression}

Specialized academic vocabulary condenses large inferential chains into single
terms. This increases efficiency within closed interpretive communities while
decreasing accessibility and obscuring the actual dependency structure of
arguments. The problem is not that compressed vocabulary exists---it is that
compressed vocabulary is often presented as if it were transparent to anyone
who has not undergone the relevant disciplinary apprenticeship.

This monograph operates under a contrary constraint. Terms are introduced
progressively, derived from previously established concepts, or explicitly
flagged as imported from identified external frameworks. Where a concept is
borrowed from sheaf theory, category theory, or field theory, the borrowing is
acknowledged and the relevant apparatus is reproduced in the appendices rather
than merely cited.

The cost is length and apparent redundancy. The benefit is that the text
functions as a reconstructable object: a reader who cannot take any prior
knowledge for granted can, in principle, reconstruct the entire conceptual
dependency graph from the text alone.

\section{The Claim-Status System as Epistemic Hygiene}

The claim-status system introduced in the Methodological Commitments chapter is
not decorative. It prevents the silent transitions from \textbf{[H]} to
\textbf{[P]} to \textbf{[D]} that produce the ontological escalation problem.

The most common failure mode in synthetic cross-domain theorizing is this: a
heuristic principle developed in one domain is applied to a second domain as a
phenomenological description; after sufficient repetition, the phenomenological
description begins functioning as a definition; the definition is then treated
as a formal result. Each transition is locally reasonable; the cumulative
effect is that a heuristic has become an ontological claim without any single
step being flagged as an escalation.

The claim-status markers interrupt this process at each step. A claim carrying
\textbf{[H]} cannot silently become a \textbf{[D]} without passing through
explicit re-evaluation. The system is an implementation of the corrigibility
condition at the level of the text itself: a mechanism for detecting and
correcting inferential drift before it propagates into downstream commitments.

\section{Memory, Continuity, and Synthetic Coherence}

The critique of persistent memory systems in AI platforms is a specific instance
of the general provenance argument. Seamless conversational continuity produced
by compressed memory summaries substitutes synthetic coherence for genuine
reconstruction. The interaction becomes smoother while the inferential structure
becomes less inspectable.

This is not a critique of memory as such but of memory that operates without
provenance transparency. A memory system that preserved the boot sequence of
prior exchanges---that made its reconstruction process visible rather than
seamless---would satisfy the provenance requirement. The concern is with systems
optimized for the appearance of continuity rather than for the reconstructibility
of the trajectory that produced it.

The analogy to narrative self-construction in biological cognition is
\textbf{[P]}: the feeling of continuous identity is also a reconstructed
coherence rather than a literal persistence. After sleep, anesthesia, distraction,
or context switching, the system reassembles a coherent identity from available
traces. The reconstruction succeeds smoothly enough that the discontinuities
become invisible. The risk, in both biological and artificial systems, is that
the reconstruction begins optimizing for smoothness rather than accuracy---that
synthetic coherence replaces grounded reconstruction.

\begin{constraintboundary}
\textbf{Established:} The connection between the monograph's writing style and
its theoretical commitments. The claim-status system as an implementation of
corrigibility at the textual level.

\textbf{Phenomenological~[P]:} The analogy between AI memory systems and
biological narrative self-construction.

\textbf{Not claimed:} That the monograph's style is the uniquely correct
implementation of these commitments. Alternative implementations that preserve
inferential continuity without progressive definition accumulation would satisfy
the same constraints.
\end{constraintboundary}
